% Generated draft: requires mathematical and source review.
\documentclass[11pt]{article}
\usepackage[margin=25mm]{geometry}
\usepackage{fontspec}
\setmainfont{Latin Modern Roman}
\usepackage{amsmath,amssymb,mathtools,unicode-math}
\setmathfont{Latin Modern Math}
\usepackage{enumitem,needspace,xurl,hyperref}
\setlength{\parindent}{0pt}
\setlength{\parskip}{5pt}
\setlength{\emergencystretch}{3em}
\newcommand{\N}{\mathbb{N}}
\newcommand{\Z}{\mathbb{Z}}
\newcommand{\Q}{\mathbb{Q}}
\newcommand{\R}{\mathbb{R}}
\newcommand{\C}{\mathbb{C}}
\newcommand{\F}{\mathbb{F}}
\newcommand{\A}{\mathbb{A}}
\newcommand{\PP}{\mathbb{P}}
\newcommand{\E}{\mathbb{E}}
\newcommand{\eps}{\varepsilon}
\newcommand{\dd}{\,\mathrm{d}}
\begin{document}
\section*{UnsolvedMath ID: 30001837}
\textbf{Source problem code:} catalogue entry
\subsection{Definitions and mathematical statement}
Let \(p\) be a prime, let \(\F_p=\Z/p\Z\), and let \(G_{\Q}=\operatorname{Gal}(\overline{\Q}/\Q)\), equipped with its profinite topology. Give \(\operatorname{GL}_2(\F_p)\) the discrete topology. A continuous representation
\[
\rho:G_{\Q}\longrightarrow\operatorname{GL}_2(\F_p)
\]
is absolutely irreducible if its scalar extension to \(\overline{\F}_p\) has no nonzero proper invariant subspace. For a complex conjugation \(c\in G_{\Q}\), call \(\rho\) even if \(\det\rho(c)=1\). This condition is independent of the choice of complex conjugation. Say that \(\rho\) is unramified outside \(p\) if \(\rho(I_\ell)=\{1\}\) for every finite prime \(\ell\ne p\), where \(I_\ell\) is an inertia subgroup at \(\ell\). No condition at \(p\) beyond continuity is specified in the supplied clean statement. Write
\[
\operatorname{SL}_2(\F_p)=\{A\in\operatorname{GL}_2(\F_p):\det A=1\}.
\]
Define the set of primes
\[
\mathcal P=\left\{p\text{ prime}:\begin{array}{l}
\text{there exists a continuous, absolutely irreducible, even}\\
\rho:G_{\Q}\to\operatorname{GL}_2(\F_p)\text{ unramified outside }p,\\
\operatorname{SL}_2(\F_p)\subseteq\rho(G_{\Q})
\end{array}\right\}.
\]
The supplied clean formulation asks the following two questions:
\begin{enumerate}
\item Is \(\mathcal P\ne\varnothing\)?
\item Is the numerical prediction
\[
\#\{p\in\mathcal P:p<10^6\}\le1
\]
true?
\end{enumerate}
The second question counts primes, not representations; it does not assert that an exceptional prime exists. The two questions are separate targets, and the numerical prediction is not a logical consequence of the existence question.

The supplied record does not restrict \(p\) to odd primes. The definition above therefore includes \(p=2\); in characteristic \(2\), the determinant condition defining evenness is automatic. Whether the intended source question excludes characteristic \(2\) requires review.

The original formulation also asks about bounds \(p\ge79\), \(k\ge33\), and unspecified lower-bound exceptions involving a Serre weight \(k\). The record does not supply the definition or normalization of \(k\), the exception list, or the hypotheses of the reported conditional bounds. These assertions cannot be reconstructed precisely from the supplied material and are not imposed on \(\mathcal P\). The supplied verification text describes conditional lower bounds for small-image cases and a surviving large-image heuristic. This description is source curation, not a mathematical verification performed for this draft.
\subsection{Short English statement}
The problem asks whether any prime \(p\) admits an absolutely irreducible continuous even two-dimensional representation of the absolute Galois group of \(\Q\), over \(\F_p\), unramified at every finite prime other than \(p\), whose image contains \(\operatorname{SL}_2(\F_p)\). It separately asks whether at most one such prime is strictly below \(10^6\). The original weight bounds and exceptions are insufficiently specified, and the intended treatment of \(p=2\) needs review. The supplied record labels the problem open and describes the numerical prediction as heuristic.
\subsection{Sources}
\begin{itemize}
\item Local UnsolvedMath registry, supplied complete record, permanent ID \(30001837\), title \emph{Large-Image Even Galois Representations}, source code \texttt{catalogue entry}, in the Oberwolfach Reports Open Problems collection. The proposed year is supplied as 2011; no individual proposer is named. The registry labels the problem \texttt{open} and the statement \texttt{corrected\_verified}; these are supplied curation labels, not independently verified conclusions.
\item Supplied source citation: \emph{Explicit Methods in Number Theory (2012)}, Oberwolfach Reports, pp. 1971--2044, DOI 10.4171/owr/2011/35. Original OWR report: \url{https://doi.org/10.4171/owr/2011/35}. The record describes Oberwolfach Reports as documenting mathematical workshops and their open-problem discussions. The citation's title year 2012 and its DOI year component 2011 are retained as supplied.
\item Supplied literature source labeled ``later partial Serre-modularity results'': \url{https://arxiv.org/abs/1908.11690}. No authors, title, precise theorem, or explanation of its applicability are supplied in the record.
\item Literature assessment checked at \texttt{2026-08-21}: this exact date identifies supplied dated curation, not an independently verified literature search. The supplied assessment labels the problem open and states that the passage records computational bounds and a heuristic conjecture about exceptional representations, rather than a proof of a global classification. The background explicitly calls this dated literature triage derived from the supplied OWR curation and requests independent verification of time-sensitive status and cited claims. Neither linked source was consulted in preparing this draft.
\end{itemize}
\end{document}
